% =============================================================================
% EBF SCOPE & STRUCTURE FRAMEWORK
% =============================================================================
%
% Meta-Framework zur Ableitung von Umfang und Struktur für jedes Thema
% 
% Für jedes Thema kann sofort abgeleitet werden:
% - Scope (Umfang) im Hauptdokument
% - Scope (Umfang) in den Appendices
% - Strukturelle Einordnung
% - Verbindungen zu anderen Komponenten
%
% =============================================================================

\documentclass[12pt,a4paper]{article}
\usepackage[margin=2.5cm]{geometry}
\usepackage{amsmath,amssymb}
\usepackage{booktabs}
\usepackage{longtable}
\usepackage{tcolorbox}
\usepackage{enumitem}
\usepackage{hyperref}
\usepackage{xcolor}

\title{EBF Scope \& Structure Framework\\
\large Meta-Framework zur Themen-Einordnung}
\author{BEATRIX Research Group}
\date{Januar 2026}

\begin{document}
\maketitle
\tableofcontents
\newpage

% =============================================================================
\section{Framework-Übersicht}
% =============================================================================

\begin{tcolorbox}[colback=blue!5!white, colframe=blue!75!black, 
    title=Zweck dieses Frameworks]

Für \textbf{jedes beliebige Thema} in EBF kann dieses Framework sofort ableiten:

\begin{enumerate}[nosep]
\item \textbf{WO} es im Dokument erscheint (Kapitel, Appendix)
\item \textbf{WIE VIEL} Raum es einnimmt (Scope-Kategorie)
\item \textbf{WELCHE STRUKTUR} es haben sollte (Template-Typ)
\item \textbf{WOMIT} es verbunden ist (Dependencies)
\end{enumerate}

\end{tcolorbox}

% =============================================================================
\section{Die EBF Architektur-Matrix}
% =============================================================================

\subsection{Ebene 1: Dokumentstruktur}

\begin{table}[htbp]
\centering
\caption{EBF Zwei-Ebenen-Architektur}
\renewcommand{\arraystretch}{1.4}
\begin{tabular}{@{}p{3cm}p{5.5cm}p{5.5cm}@{}}
\toprule
\textbf{Eigenschaft} & \textbf{Hauptdokument} & \textbf{Appendices} \\
\midrule
Zweck & Konzeptuelle Erzählung & Technische Spezifikation \\
Stil & Zugänglich, narrativ & Formal, erschöpfend \\
Leser & Ökonomen, Policy Maker & Spezialisten, Theoretiker \\
Tiefe & Intuition + Kerngleichungen & Alle Axiome + Beweise \\
Umfang & ~485 Seiten (70 Dateien) & ~750 Seiten (56 Dateien) \\
\bottomrule
\end{tabular}
\end{table}

\subsection{Ebene 2: Themen-Hierarchie}

\begin{tcolorbox}[colback=gray!5!white, colframe=gray!75!black, 
    title=Die 4 fundamentalen Fragen]

Jedes Thema in EBF adressiert mindestens eine dieser Fragen:

\begin{center}
\renewcommand{\arraystretch}{1.5}
\begin{tabular}{@{}clll@{}}
\toprule
\textbf{Code} & \textbf{Frage} & \textbf{Antwort} & \textbf{Symbol} \\
\midrule
WHO & Wer hat Utility? & Aggregationsebenen & $L \in \{0,1,...,\infty\}$ \\
WHAT & Was ist Utility? & FEPSDE-Dimensionen & $d \in \{F,E,P,S,D,E\}$ \\
HOW & Wie interagieren sie? & Complementarity & $\gamma_{ij}$ \\
WHEN & Wann wirkt Kontext? & Ψ-Dimensionen & $\Psi = (\Psi_1,...,\Psi_8)$ \\
\bottomrule
\end{tabular}
\end{center}

\end{tcolorbox}

% =============================================================================
\section{Scope-Kategorien (Umfang)}
% =============================================================================

\subsection{Hauptdokument: 6 Scope-Kategorien}

\begin{table}[htbp]
\centering
\caption{Scope-Kategorien im Hauptdokument}
\renewcommand{\arraystretch}{1.3}
\begin{tabular}{@{}clccp{5cm}@{}}
\toprule
\textbf{Kat.} & \textbf{Name} & \textbf{Seiten} & \textbf{Dateien} & \textbf{Beispiele} \\
\midrule
S1 & Paragraph & 0.5--2 & 0 (inline) & Einzelne Definition, Beispiel \\
S2 & Subsection & 2--10 & 0--1 & Einzelnes Konzept \\
S3 & Section & 10--25 & 1 & Ganzes Kapitel (z.B. Ch. 5) \\
S4 & Multi-Section & 25--50 & 2--5 & Kapitel mit Subkapiteln \\
S5 & Major Chapter & 50--150 & 10--20 & Ch. 10 (Welfare), Ch. 11 (Awareness) \\
S6 & Part & 150+ & 20+ & Teil I--VI \\
\bottomrule
\end{tabular}
\end{table}

\subsection{Appendices: 5 Scope-Kategorien}

\begin{table}[htbp]
\centering
\caption{Scope-Kategorien in Appendices}
\renewcommand{\arraystretch}{1.3}
\begin{tabular}{@{}clccp{5cm}@{}}
\toprule
\textbf{Kat.} & \textbf{Name} & \textbf{Seiten} & \textbf{KB} & \textbf{Beispiele} \\
\midrule
A1 & Stub & 1--5 & <5 & D (Proofs), E (Ops) - unvollständig \\
A2 & Basic & 5--20 & 5--20 & H (History), T (Meta) \\
A3 & Standard & 20--50 & 20--50 & Meiste DOMAIN, LIT Appendices \\
A4 & Extended & 50--100 & 50--100 & AAA (WHO), C (WHAT) \\
A5 & Comprehensive & 100+ & 150+ & AU (Axioms), AV (Willingness) \\
\bottomrule
\end{tabular}
\end{table}

% =============================================================================
\section{Struktur-Templates}
% =============================================================================

\subsection{Hauptdokument: Chapter-Struktur}

\begin{tcolorbox}[colback=green!5!white, colframe=green!75!black, 
    title=Standard Chapter Template]

\textbf{OPENING} (Pflicht)
\begin{itemize}[nosep]
\item Section heading + label
\item Opening paragraph (WAS und WARUM)
\item Chapter roadmap
\end{itemize}

\textbf{CORE CONTENT}
\begin{itemize}[nosep]
\item Subsections für Hauptthemen
\item Pattern: Intuition → Beispiel → Formal → Interpretation
\item Kerngleichungen (nummeriert, interpretiert)
\item Verbindungen zu anderen Kapiteln
\end{itemize}

\textbf{CLOSING} (Pflicht)
\begin{itemize}[nosep]
\item Summary
\item Appendix-Verweise
\item Überleitung zum nächsten Kapitel
\end{itemize}

\end{tcolorbox}

\subsection{Appendices: 4-Teil-Struktur}

\begin{tcolorbox}[colback=yellow!5!white, colframe=yellow!75!black, 
    title=Standard Appendix Template]

\textbf{PART 1: FRONT MATTER} (Pflicht)
\begin{itemize}[nosep]
\item Header Block (Metadata)
\item Chapter Linkage Box
\item Cross-Reference Map
\item Abstract (max 100 Wörter)
\item Quick Reference Box
\end{itemize}

\textbf{PART 2: CORE CONTENT} (Appendix-spezifisch)
\begin{itemize}[nosep]
\item Fundamental Question
\item Core Theory + Axiome
\item Key Results
\item Integration mit anderen Appendices
\end{itemize}

\textbf{PART 3: APPLICATION} (Pflicht)
\begin{itemize}[nosep]
\item Worked Example
\item Implications for Practice
\end{itemize}

\textbf{PART 4: BACK MATTER} (Pflicht)
\begin{itemize}[nosep]
\item Summary Box
\item Glossary of Symbols
\item Critical Foundations (für CORE)
\item Open Issues
\item References
\end{itemize}

\end{tcolorbox}

% =============================================================================
\section{Themen-Klassifikation: Die 8 Appendix-Kategorien}
% =============================================================================

\begin{table}[htbp]
\centering
\caption{Die 8 Appendix-Kategorien mit Eigenschaften}
\renewcommand{\arraystretch}{1.4}
\small
\begin{tabular}{@{}clcccp{4cm}@{}}
\toprule
\textbf{Kat.} & \textbf{Name} & \textbf{Anzahl} & \textbf{Scope} & \textbf{Pflicht} & \textbf{Zweck} \\
\midrule
CORE & Fundamentaltheorie & 4 & A4--A5 & 7 Req. & Beantwortet WHO/WHAT/HOW/WHEN \\
FORMAL & Mathematik & 4 & A2--A5 & Beweise & Formale Ableitungen \\
DOMAIN & Anwendung & 13 & A3 & Standard & Feldspezifische Anwendung \\
CONTEXT & Ψ-Framework & 3 & A3--A4 & Standard & Zeit/Raum/Kontext \\
METHOD & Methodik & 4 & A2--A3 & Standard & Empirische Methoden \\
PREDICT & Vorhersagen & 7 & A3 & Falsifizierbar & Testbare Predictions \\
LIT & Literatur & 10 & A2--A3 & Vollständig & Autor/Themen-Integration \\
REF & Nachschlagen & 4 & A2--A3 & Standard & Glossar, Beispiele \\
\bottomrule
\end{tabular}
\end{table}

% =============================================================================
\section{Die Scope-Ableitung: Entscheidungsbaum}
% =============================================================================

\begin{tcolorbox}[colback=red!5!white, colframe=red!75!black, 
    title=Entscheidungsbaum: Wo gehört ein Thema hin?]

\textbf{Schritt 1: Welche fundamentale Frage?}
\begin{itemize}[nosep]
\item WHO → Levels, Aggregation → Ch. 10.8 + CORE-WHO (AAA)
\item WHAT → Dimensionen, Utility → Ch. 10.1--10.7 + CORE-WHAT (C)
\item HOW → Interaktion, Complementarity → Ch. 5 + CORE-HOW (B)
\item WHEN → Kontext, Ψ → Ch. 9 + CORE-WHEN (V)
\end{itemize}

\textbf{Schritt 2: Welche Tiefe?}
\begin{itemize}[nosep]
\item Konzeptuell → Hauptdokument (entsprechendes Kapitel)
\item Formal/Vollständig → Appendix (entsprechende Kategorie)
\end{itemize}

\textbf{Schritt 3: Welcher Typ?}
\begin{itemize}[nosep]
\item Neue Theorie → CORE oder FORMAL
\item Anwendungsgebiet → DOMAIN
\item Methodik → METHOD
\item Testbare Hypothese → PREDICT
\item Literaturintegration → LIT
\end{itemize}

\end{tcolorbox}

% =============================================================================
\section{Schnellreferenz: Thema → Struktur}
% =============================================================================

\subsection{Matrix: Thementyp × Scope}

\begin{longtable}{@{}p{4cm}ccp{3cm}p{3cm}@{}}
\caption{Thementyp-Scope-Matrix} \\
\toprule
\textbf{Thementyp} & \textbf{Ch. Scope} & \textbf{App. Scope} & \textbf{Chapter} & \textbf{Appendix} \\
\midrule
\endfirsthead
\midrule
\textbf{Thementyp} & \textbf{Ch. Scope} & \textbf{App. Scope} & \textbf{Chapter} & \textbf{Appendix} \\
\midrule
\endhead
\bottomrule
\endlastfoot

% FUNDAMENTALE KONZEPTE
\multicolumn{5}{@{}l}{\textit{\textbf{Fundamentale Konzepte}}} \\
Aggregationsebenen (L) & S4 & A4 & 10.8 & CORE-WHO \\
Utility-Dimensionen (FEPSDE) & S5 & A4 & 10 & CORE-WHAT \\
Complementarity ($\gamma$) & S3 & A3→A4 & 5 & CORE-HOW \\
Kontext ($\Psi$) & S4 & A3→A4 & 9 & CORE-WHEN \\
\addlinespace

% UTILITY-KOMPONENTEN
\multicolumn{5}{@{}l}{\textit{\textbf{Utility-Komponenten}}} \\
Individual Utility (INU) & S2 & A3 & 10.1 & C (Section) \\
Collective Utility (KNU) & S2 & A3 & 10.1 & C (Section) \\
Identity Utility (IDN) & S2 & A3 & 10.1 & C (Section) \\
Financial (F) & S1 & A2 & 10.2 & C (Section) \\
Emotional (E) & S1 & A2 & 10.2 & C (Section) \\
Physical (P) & S1 & A2 & 10.2 & C (Section) \\
Social (S) & S1 & A2 & 10.2 & C (Section) \\
Digital (D) & S1 & A2 & 10.2 & C (Section) \\
Ecological (E) & S1 & A2 & 10.2 & C (Section) \\
\addlinespace

% MATHEMATISCHE KONZEPTE
\multicolumn{5}{@{}l}{\textit{\textbf{Mathematische Konzepte}}} \\
Supermodularität & S2 & A2 & 5, 8 & A, X \\
Referenzpunkte & S3 & A2 & 6 & A \\
Kohärenz-Index (K) & S2 & A2 & 8 & A \\
Qualitäts-Index (Q) & S2 & A2 & 10 & A \\
\addlinespace

% KONTEXT-DIMENSIONEN
\multicolumn{5}{@{}l}{\textit{\textbf{Kontext-Dimensionen ($\Psi_1$--$\Psi_8$)}}} \\
Cognitive Load ($\Psi_1$) & S1 & A2 & 9 & V (Section) \\
Time Pressure ($\Psi_2$) & S1 & A2 & 9 & V (Section) \\
Social Observation ($\Psi_3$) & S1 & A2 & 9 & V (Section) \\
Resource Scarcity ($\Psi_4$) & S1 & A2 & 9 & V (Section) \\
Emotional State ($\Psi_5$) & S1 & A2 & 9 & V (Section) \\
Physical Environment ($\Psi_6$) & S1 & A2 & 9 & V (Section) \\
Information Availability ($\Psi_7$) & S1 & A2 & 9 & V (Section) \\
Cultural Frame ($\Psi_8$) & S1 & A2 & 9 & V (Section) \\
\addlinespace

% AWARENESS/WILLINGNESS
\multicolumn{5}{@{}l}{\textit{\textbf{Awareness \& Willingness}}} \\
Awareness Function (A) & S5 & A5 & 11 & AV \\
Willingness Threshold (W) & S4 & A5 & 12 & AV \\
Motivated Beliefs & S3 & A3 & 11.03 & AV (Section) \\
Biases \& Distortions & S2 & A2 & 11.11 & AV (Section) \\
\addlinespace

% ANWENDUNGSGEBIETE
\multicolumn{5}{@{}l}{\textit{\textbf{Anwendungsgebiete (DOMAIN)}}} \\
Labor Economics & S2 & A3 & 10.9.1 & DOMAIN-LABOR \\
Matching Markets & S2 & A3 & 5 & DOMAIN-MATCH \\
Industrial Organization & S2 & A3 & 5 & DOMAIN-IO \\
Game Theory & S2 & A3 & 5 & DOMAIN-EVO \\
Mechanism Design & S2 & A3 & 5 & DOMAIN-MECH \\
Social Choice & S2 & A3 & 10.8 & DOMAIN-CHOICE \\
Health & S2 & A2 & 10.9.2 & REF-EXAMPLES \\
Finance & S2 & A2 & 10.9.3 & REF-EXAMPLES \\
Sustainability & S2 & A2 & 10.9.4 & REF-EXAMPLES \\
\addlinespace

% METHODIK
\multicolumn{5}{@{}l}{\textit{\textbf{Methodik (METHOD)}}} \\
LLM Monte Carlo & S3 & A3 & 4x & METHOD-LLMMC \\
Operationalisierung & S2 & A2 & 4 & METHOD-OPS \\
Evaluation Protocol & S2 & A3 & 4 & METHOD-EVAL \\
Self-Reinforcement Learning & S2 & A3 & 11 & METHOD-SRL \\
\addlinespace

% LITERATUR
\multicolumn{5}{@{}l}{\textit{\textbf{Literatur-Integration (LIT)}}} \\
Nobel Contributions & S1 & A3 & 1 & LIT-NOBEL \\
Fehr Papers & S1 & A3 & 4 & LIT-FEHR \\
Acemoglu Papers & S1 & A3 & 9 & LIT-ACEMOGLU \\
Kahneman/Tversky & S1 & A3 & 3, 6 & LIT-KT \\
\addlinespace

% PREDICTIONS
\multicolumn{5}{@{}l}{\textit{\textbf{Predictions (PREDICT)}}} \\
Trade War & S2 & A3 & 11.1 & PREDICT-01 \\
Spillovers & S2 & A3 & 11.2 & PREDICT-02 \\
Asymmetric Shocks & S2 & A3 & 11.3 & PREDICT-03 \\
Techlash & S2 & A3 & 11.4 & PREDICT-04 \\
Identity Polarization & S2 & A3 & 11.5 & PREDICT-05 \\
Coherence Trap & S2 & A3 & 11.6 & PREDICT-06 \\

\end{longtable}

% =============================================================================
\section{CORE Appendix: Spezielle Anforderungen}
% =============================================================================

\begin{tcolorbox}[colback=red!5!white, colframe=red!75!black, 
    title=Die 7 CORE-Anforderungen (zusätzlich zum Standard)]

Ein Appendix wird CORE wenn er \textbf{ALLE 7} erfüllt:

\begin{enumerate}
\item \textbf{Fundamentale Frage:} Beantwortet WHO, WHAT, HOW, oder WHEN
\item \textbf{Axiom-System:} Eigene Axiome (8--200+) mit Codes (AL-, γ-, Ψ-)
\item \textbf{Literatur:} 80--150+ Referenzen, thematisch organisiert
\item \textbf{Critical Foundations:} 5--10 steelmanned Einwände + Antworten
\item \textbf{Open Issues:} 5--8 offene Fragen mit Roadmap
\item \textbf{Bidirektionale Integration:} Verbindung zu ALLEN anderen CORE
\item \textbf{Glossar:} 20+ Symboldefinitionen
\end{enumerate}

\textbf{Aktueller Status:}
\begin{center}
\begin{tabular}{@{}lcccc@{}}
\toprule
Requirement & WHO (AAA) & WHAT (C) & HOW (B) & WHEN (V) \\
\midrule
Fundamentale Frage & ✓ & ✓ & ✓ & ✓ \\
Axiom-System & ✓ (8) & ✓ (220) & ✗ & ✗ \\
Literatur & ✓ (483) & ✓ (170) & ✗ & ✗ \\
Critical Foundations & ✓ (10) & ✗ & ✗ & ✗ \\
Open Issues & ✓ (8) & ✗ & ✗ & ✗ \\
Integration & ~50\% & ~50\% & ✗ & ✗ \\
Glossar & ✓ (29) & ✓ (50) & ✗ & ✗ \\
\midrule
\textbf{Erfüllung} & \textbf{70\%} & \textbf{45\%} & \textbf{10\%} & \textbf{15\%} \\
\bottomrule
\end{tabular}
\end{center}

\end{tcolorbox}

% =============================================================================
\section{Praktische Anwendung: Beispiele}
% =============================================================================

\subsection{Beispiel 1: Neues Thema ``Intergenerationale Transfers''}

\begin{tcolorbox}[colback=blue!5!white, colframe=blue!50!black]
\textbf{Frage:} Wo gehört ``Intergenerationale Transfers'' hin?

\textbf{Analyse:}
\begin{enumerate}[nosep]
\item Fundamentale Frage: WHO (Ebene $L=\infty$ Meta) + WHAT (F-Dimension)
\item Tiefe: Anwendungsgebiet, nicht neue Theorie
\item Typ: DOMAIN
\end{enumerate}

\textbf{Ergebnis:}
\begin{itemize}[nosep]
\item Hauptdokument: Ch. 10.8 (Aggregate Welfare), Scope S2
\item Appendix: DOMAIN-Kategorie, neuer Appendix oder Section in AAA
\item Scope: A2--A3 (10--30 Seiten)
\item Verbindungen: CORE-WHO (Ebenen), CORE-WHAT (F-Dimension)
\end{itemize}
\end{tcolorbox}

\subsection{Beispiel 2: Neues Thema ``Visceral Factors''}

\begin{tcolorbox}[colback=blue!5!white, colframe=blue!50!black]
\textbf{Frage:} Wo gehört ``Visceral Factors'' hin?

\textbf{Analyse:}
\begin{enumerate}[nosep]
\item Fundamentale Frage: WHEN (Kontext $\Psi_5$ Emotional State)
\item Tiefe: Erweiterung der Kerntheorie
\item Typ: CORE-Erweiterung oder CONTEXT
\end{enumerate}

\textbf{Ergebnis:}
\begin{itemize}[nosep]
\item Hauptdokument: Ch. 9 (Context), Scope S2--S3
\item Appendix: CORE-WHEN (V), erweiterte Section für $\Psi_5$
\item Scope: A2 (5--15 Seiten innerhalb V)
\item Verbindungen: CORE-WHEN, Ch. 11 (Awareness)
\end{itemize}
\end{tcolorbox}

\subsection{Beispiel 3: Neues Thema ``Digital Platform Economics''}

\begin{tcolorbox}[colback=blue!5!white, colframe=blue!50!black]
\textbf{Frage:} Wo gehört ``Digital Platform Economics'' hin?

\textbf{Analyse:}
\begin{enumerate}[nosep]
\item Fundamentale Frage: HOW (Complementarity in Networks) + WHAT (D-Dimension)
\item Tiefe: Anwendungsgebiet mit eigenem Appendix
\item Typ: DOMAIN
\end{enumerate}

\textbf{Ergebnis:}
\begin{itemize}[nosep]
\item Hauptdokument: Ch. 5 (Complementarity), Ch. 10.2 (D-Dimension), Scope S2
\item Appendix: Neuer DOMAIN-DIGITAL Appendix
\item Scope: A3 (20--40 Seiten)
\item Verbindungen: CORE-HOW, CORE-WHAT (D), DOMAIN-IO
\end{itemize}
\end{tcolorbox}

% =============================================================================
\section{Quick Reference Cards}
% =============================================================================

\subsection{Card 1: Scope-Bestimmung}

\begin{tcolorbox}[colback=gray!10!white, colframe=black]
\textbf{SCOPE QUICK REFERENCE}

\textbf{Hauptdokument:}
\begin{itemize}[nosep]
\item S1 (0.5--2 S.): Einzelne Definition
\item S2 (2--10 S.): Ein Konzept
\item S3 (10--25 S.): Ein Kapitel
\item S4 (25--50 S.): Kapitel + Subkapitel
\item S5 (50--150 S.): Major Chapter
\item S6 (150+ S.): Part
\end{itemize}

\textbf{Appendix:}
\begin{itemize}[nosep]
\item A1 (<5 S.): Stub
\item A2 (5--20 S.): Basic
\item A3 (20--50 S.): Standard
\item A4 (50--100 S.): Extended
\item A5 (100+ S.): Comprehensive
\end{itemize}
\end{tcolorbox}

\subsection{Card 2: Kategorien-Bestimmung}

\begin{tcolorbox}[colback=gray!10!white, colframe=black]
\textbf{KATEGORIE QUICK REFERENCE}

\textbf{Ist es fundamental für EBF?}
\begin{itemize}[nosep]
\item Ja → CORE (WHO/WHAT/HOW/WHEN)
\item Nein → Weiter
\end{itemize}

\textbf{Braucht es mathematische Beweise?}
\begin{itemize}[nosep]
\item Ja → FORMAL
\item Nein → Weiter
\end{itemize}

\textbf{Ist es ein Anwendungsgebiet?}
\begin{itemize}[nosep]
\item Ja → DOMAIN
\item Nein → Weiter
\end{itemize}

\textbf{Betrifft es Ψ-Dimensionen?}
\begin{itemize}[nosep]
\item Ja → CONTEXT
\item Nein → Weiter
\end{itemize}

\textbf{Ist es eine Methodik?}
\begin{itemize}[nosep]
\item Ja → METHOD
\item Nein → Weiter
\end{itemize}

\textbf{Ist es testbar/falsifizierbar?}
\begin{itemize}[nosep]
\item Ja → PREDICT
\item Nein → Weiter
\end{itemize}

\textbf{Integriert es Literatur?}
\begin{itemize}[nosep]
\item Ja → LIT
\item Nein → REF
\end{itemize}
\end{tcolorbox}

\subsection{Card 3: Struktur-Bestimmung}

\begin{tcolorbox}[colback=gray!10!white, colframe=black]
\textbf{STRUKTUR QUICK REFERENCE}

\textbf{Hauptdokument-Kapitel:}
\begin{enumerate}[nosep]
\item OPENING: Section + Opening + Roadmap
\item CONTENT: Intuition → Beispiel → Formal → Interpretation
\item CLOSING: Summary + Appendix-Refs + Überleitung
\end{enumerate}

\textbf{Standard-Appendix:}
\begin{enumerate}[nosep]
\item FRONT: Header + Chapter-Link + Abstract + Quick-Ref
\item CORE: Question + Theory + Axioms + Results + Integration
\item APPLICATION: Example + Implications
\item BACK: Summary + Glossar + Issues + Refs
\end{enumerate}

\textbf{CORE-Appendix (zusätzlich):}
\begin{itemize}[nosep]
\item 7 CORE-Anforderungen erfüllen
\item Bidirektionale Links zu allen anderen CORE
\item 80--150+ Referenzen
\item Critical Foundations (5--10 Objections)
\end{itemize}
\end{tcolorbox}

% =============================================================================
\section{Anhang: Vollständige Themen-Matrix}
% =============================================================================

% Dies wird eine umfassende Tabelle aller ~100+ Themen mit ihrer Einordnung
% [Hier würde die vollständige Matrix eingefügt]

\end{document}
